\chapter{Sistemas de ecuaciones lineales}
\label{cap:sistemas}

La eliminación de Gauss se interpreta ahora como un método para estudiar aplicaciones lineales y conjuntos de soluciones. Así se conectan los algoritmos matriciales con núcleo, imagen y rango \cite{lang1986,hoffmankunze1971}.

\section{Sistemas lineales y ecuaciones matriciales}
Representación de sistemas como $Ax=b$ y lectura por columnas.

\section{Operaciones elementales}
Operaciones por renglones y su efecto sobre las soluciones.

\section{Eliminación gaussiana}
Pivotes, variables libres y reducción de sistemas.

\section{Forma escalonada y forma escalonada reducida}
Definiciones, existencia y uso para describir el conjunto solución.

\section{Sistemas homogéneos}
Espacio de soluciones, núcleo de una matriz y variables libres.

\section{Sistemas no homogéneos}
Compatibilidad y estructura afín del conjunto solución.

\section{Existencia y unicidad de soluciones}
Criterios de compatibilidad y unicidad expresados mediante rango y nulidad.

\section{Relación entre \texorpdfstring{$Ax=b$}{Ax=b}, núcleo e imagen}
Existencia cuando $b$ pertenece a la imagen y descripción de todas las soluciones a partir del núcleo.
